
\begin{abstract}
本文提出一套以 Zeckendorf 截断折叠为规范（canonical folding）的时空结构图接口：把微观态空间 $X$、宏观观测空间 $Z$、折叠算子 $f:X\to Z$、空间关系与时间关系统一编码为一个带颜色的组合对象，使其可直接接入 Babai 的 \WL{} / 相干配置框架。该接口将“观测导致的信息压缩”表述为由 $f$ 诱导的块系统，并将纤维 $f^{-1}(z)$ 的内部自由度解释为可计算的“时间/复杂度”坐标。\par
在规范体系下，我们给出结构同构与编码图同构的等价翻译定理，并给出可复现实验指标：\WL{} 稳定深度、残余轨道规模与纤维统计量，从而为“带时空语义的离散对象”接入 Babai--\WL{}--相干配置框架提供一个可审计的最小接口。
\end{abstract}

\begin{sloppypar}
\textbf{关键词：} 时空结构图；观测折叠；Zeckendorf 截断；图同构；Weisfeiler--Leman；相干配置；残余对称；可复现实验
\end{sloppypar}

\paragraph{约定与记号。}
除非另行说明，$\log$ 表示自然对数。微观二进制串采用下标从 $0$ 开始的位序 $x=x_0x_1\cdots x_{m-1}$，并定义整数解释
$$
N(x)=\sum_{i=0}^{m-1} x_i\,2^i.
$$
本文的 Fibonacci 数列采用 Zeckendorf 版本：$F_1=1$，$F_2=2$，$F_{k+2}=F_{k+1}+F_k$；Zeckendorf 数码记为 $\mathbf c(N)=(c_1(N),c_2(N),\dots)$。关于 \GI{}、\WL{}、相干配置与 Babai 框架的背景参见 \cite{babai2016gi,wl2018symmetryvsregularity,ivanov2013coherent}；Zeckendorf 表示参见 \cite{zeckendorf1972representation}。

